\documentclass[12pt,a4paper]{article}

\usepackage{geometry}
\geometry{
    a4paper,
    left=25mm,
    right=25mm,
}
\nocite{*}


\linespread{1.3}


% Packages
\usepackage{amsmath}
\usepackage{tikz-cd}
\usepackage{amsfonts}
\usepackage{amsthm}
\usepackage{amssymb}
\usepackage{graphicx}
\usepackage{threeparttable}
\usepackage{caption}
\usepackage{booktabs}
\usepackage{xcolor}
\usepackage{multirow}
\usepackage{hyperref}
\hypersetup{colorlinks=true, citecolor=blue}

\usepackage{biblatex}
\usepackage{textcomp}
\usepackage{mathtools}
\usepackage{stmaryrd}
\addbibresource{bib.bib}

\newenvironment{itemizeb}
{\begin{itemize}
     \itemsep=2pt}{
\end{itemize}}

\numberwithin{equation}{section}


\newtheoremstyle{mytheorem}
{}% measure of space to leave above the theorem. E.g.: 3pt
{}% measure of space to leave below the theorem. E.g.: 3pt
{\it}% name of font to use in the body of the theorem
{\parindent}% measure of space to indent
{\bf}% name of head font
{.}% punctuation between head and body
{ }% space after theorem head; " " = normal interword space
{\thmnumber{#2.~}\thmname{#1}\thmnote{~\rm#3}}% Manually specify head

\newtheoremstyle{myremark}
{}% measure of space to leave above the theorem. E.g.: 3pt
{}% measure of space to leave below the theorem. E.g.: 3ptontinuo dentro e fuori dall'infermeria, poi Barzagli e i due mesi di stop per l'infortunio alla spalla, infine a
{\rm}% name of font to use in the body of the theorem
{\parindent}% measure of space to indent
{\bf}% name of head font
{.}% punctuation between head and body
{ }% space after theorem head; " " = normal interword space
{\thmnumber{#2.~}\thmname{#1}\thmnote{~\rm#3}}% Manually specify head

\newtheoremstyle{myparagraph}
{}% measure of space to leave above the theorem. E.g.: 3pt
{}% measure of space to leave below the theorem. E.g.: 3pt
{\rm}% name of font to use in the body of the theorem
{\parindent}% measure of space to indent
{\bf}% name of head font
{.}% punctuation between head and body
{ }% space after theorem head; " " = normal interword space
{\thmnumber{#2.~}\thmname{#1}\thmnote{#3}}% Manually specify head

\theoremstyle{mytheorem}
\newtheorem{maintheorem}{Theorem}
\newtheorem{theorem}[subsection]{Theorem}
\newtheorem{definition}[subsubsection]{Definition}
\newtheorem{lemma}[subsection]{Lemma}
\newtheorem{corollary}[subsection]{Corollary}
\newtheorem{proposition}[subsection]{Proposition}
\newtheorem{question}[subsection]{Question}

\theoremstyle{myremark}
\newtheorem{remark}[subsection]{Remark}
%\newtheorem*{remark*}{Remark}
%\newtheorem{remarks}[subsection]{Remarks}
%\newtheorem*{remarks*}{Remarks}
\newtheorem{problem}[subsection]{Problem}
\theoremstyle{myparagraph}
\newtheorem{parag}[subsection]{}
\newtheorem*{parag*}{}


\renewcommand\qedsymbol{$\boxempty$}


\newcommand{\R}{\mathbb{R}}

\usepackage{pifont} % For tick and cross symbols
% Define commands for tick and cross
\newcommand{\yes}{\ding{51}}
\newcommand{\no}{\ding{55}}
\newcommand{\na}{---} % For not applicable



\author{Dario Filatrella}

\title{Branched optimal transport}
\date{August 2024}

\begin{document}

    \large

    \pagenumbering{gobble}

    \maketitle

    \newpage

    \tableofcontents

    \newpage

    \pagenumbering{arabic}
    \setcounter{page}{1}
    \thispagestyle{empty}


    \section{Note burocratiche per la Bocconi}

    un lavoro (tesi di ricerca) atto a produrre nuova conoscenza o nuove metodologie
    scientifiche, o finalizzato ad analizzare un problema e a fornirne adeguata soluzione.
    Il lavoro deve includere una approfondita review della letteratura, la definizione di un
    problema di ricerca basato su un research gap, la metodologia di ricerca utilizzata,
    l’analisi dei risultati ottenuti.
    La tesi si sviluppa orientativamente in 50 pagine (circa 18.000 parole).

    Durante la seduta di laurea il candidato può distribuire materiali alla commissione (indicativamente un
    massimo di 12 slides comprensive di tabelle, grafici, etc.) che possano rappresentare un supporto utile
    alla discussione. A questo proposito, si raccomanda di utilizzare solo materiale cartaceo e riciclabile
    (evitando copertine e rilegature plastificate). Si ricorda che non è consentita la proiezione di slides
    durante la presentazione.

    Cv@b da compilare (opzionale)

    inserisce l’abstract che sintetizza l’oggetto del lavoro in un massimo di 4000
    caratteri - vedi NOTA IMPORTANTE;

    Il file definitivo, in formato pdf, deve contenere, in testa, 4 pagine così ordinate
    I pagina: una pagina bianca;
    II pagina una pagina bianca;
    III pagina: dedica o ringraziamenti (oppure una pagina bianca);
    IV pagina: una pagina bianca.
    Il testo del proprio lavoro (indice, premessa, contenuto etc) DEVE iniziare, numerazione delle
    pagine compresa, DOPO le 4 pagine sopradescritte; inoltre, il contenuto della tesi non deve
    includere il frontespizio che verrà riprodotto automaticamente in fase di stampa e riporterà le
    informazioni sul titolo inserite dallo studente online, nè l'abstract.
    Inoltre, deve rispettare i seguenti parametri:
    Dimensione massima del file: 10 MB
    Formato: cm 29x21 (A4)
    Margine destro e sinistro 2.5 cm
    Battitura pagine: da 26 a 30 righe
    Carattere consigliato Arial/Tahoma/Verdana
    Corpo 12
    Pagine numerate


    \section{Background}
    \begin{itemize}
        \item discorso organico sull-argomento
        \item introduzione storica potenzialmente
        \item obbiettivo del lavoro (stabilità dei risultati e lato computazionale)
    \end{itemize}

    TODO: FA DAVVERO PENA
    The first mathematical formulation of an \textit{optimal transport} problem was given by Gaspard Monge in 1781~\cite[Monge]{monge1781mémoire}.
    In his work, Monge considered the problem of moving a pile of soil from one location to another (in fact, a special case of the problem is known as the \textit{Earth-mover distance}).
    The modern formulation of the problem is due to Kantorovich~\cite[Kantorovich 1958]{Kantorovich1958}, almost two centuries later.
    The language used by Kantorovich is that of \textit{measure theory} and formally the problem is stated as follows:
    given $X \subset \R^d$, two probability measures $\mu_0, \mu_1$ on $X$ with finite and equal mass, a cost function $c: X \times X \to \R^+$ find a transport map $T: X \to X$ that realizes the following infimum:
    \begin{equation}
        \inf \left\{ \int_X c(x, T(x)) d\mu_0(x) \, : \, T_\# \mu_0 = \mu_1 \right\}
        \label{eq:monge}
    \end{equation}

    Intuitively the cost function $c$ represents the cost of moving a unit mass from $x$ to $y$, and the functional to be minimized represents is the sum (integral)
    over all the particles that need to be moved.
    There are of course many variations, we could add assumptions on $X$ (compactness, convexity, etc.) or remove assumptions (considering $X$ to be a generic metric space).


    {\color{red}TODO: la bibliografia non funziona non compare Monge}


    \section{Literature review}


    \section{More advanced results (stability)}


    \section{Literature review computazionale e numerica}


    \section{Miei risultati teorici (con anche controesempi per spiegare le difficoltà numeriche)}


    \section{Risultati numerici}


    \section{Conclusioni e future development}


    \newpage

    \printbibliography


\end{document}